\paragraph{Survivor procedure.}
Table~\ref{tab:algorithm} gives the implemented selection procedure. The gate is
sequential: an already removed candidate is skipped as a potential dominator.
This ordering prevents a deleted prompt from supplying a later elimination
certificate, but it also means the procedure is not an order-invariant batch
solution to all possible pairwise inequalities. Stable prompt IDs make the
chosen order reproducible. They do not provide an optimality guarantee.

\begin{table}[tbp]\centering\small
\caption{PACE survivor procedure. This is pseudocode for one invocation, not the
complete evaluation scheduler. $\operatorname{NSGA}(R,k)$ denotes rank/crowding
selection of at most $k$ prompt IDs from $R$.}
\label{tab:algorithm}
\begin{tabularx}{\linewidth}{@{}r>{\raggedright\arraybackslash}X@{}}\toprule
Step & Operation\\\midrule
1 & Input candidate IDs $C$, observed state, active capacity $\mu$, archive allowance
$K$, threshold $\varepsilon$, and seeded bootstrap generator. Set $E=\varnothing$.\\
2 & Visit $q$ in sorted $C$. Visit $p$ in sorted $C$, skipping $p=q$ and IDs in $E$.\\
3 & Skip the pair if $\widehat c(p)>\widehat c(q)$, if no common samples exist,
or if fewer than eight common samples are available.\\
4 & If $L_{pq}>\varepsilon$, append $(q,p,L_{pq},U_{pq},n)$ to the confidence log,
set $E\leftarrow E\cup\{q\}$, and stop scanning dominators for $q$.\\
5 & Set $R=C\setminus E$. If $|R|>\mu+K$, keep $\operatorname{NSGA}(R,\mu+K)$
and log all other IDs as forced evictions.\\
6 & If $|R|>\mu$, set $A=\operatorname{NSGA}(R,\mu)$; otherwise set $A=R$.
Return retained set $R$, active set $A$, and the two removal logs.\\\bottomrule
\end{tabularx}
\end{table}

The control B0+Archive executes the same capacity and active-set steps with
$E=\varnothing$. In contrast, archive-free B0 directly truncates to $\mu$.
Consequently, disabling the gate while retaining $K=8$ is the relevant primary
control, whereas comparing PACE against archive-free B0 changes two components.
Setting $K=0$ removes headroom but does not logically disable evidence-based
elimination; the gate can still change which candidates reach NSGA-II truncation.
Any similarity under that ablation is an observation, not an identity of the rules.

An invocation can return fewer than $\mu$ active prompts if confidence elimination
leaves too few candidates. The implementation does not repopulate the set with
eliminated IDs to fill empty slots. This detail matters when interpreting
population size as an upper bound rather than an invariant count at every phase.
Likewise, a candidate retained in the archive is not automatically selected as a
parent: only the active set enters the variation operator.

\subsection{Variation operators and representation constraints}
Each instruction contains a literal \texttt{\{text\}} placeholder, which is
replaced with the article at evaluation time. The online initial population is
\texttt{p00}, \texttt{p01}, \texttt{p02}, \texttt{p03}, \texttt{p04},
\texttt{p05}, \texttt{p07}, and \texttt{p15}. This eight-prompt population is a
subset of the 16-prompt pool used in offline replay; the two pools should not be
conflated when assessing what candidates were available.

For each child, two parent IDs are sampled with replacement from the active
population. With probability one half, crossover combines an instruction head
from one parent with a response-format tail from the other, splitting at the
placeholder. Otherwise, the first parent's text is copied. Every child then
receives one of five uniformly chosen mutation types: replace or add a role
prefix; alter the response-format suffix; add a label-definition sentence;
make a small wording substitution; or shorten the instruction using line-based
chunks. Placeholder repair appends an article field when needed.

These operations generate readable strings without extra LLM calls. They are
intended to preserve the task and label vocabulary, but they do not prove semantic
equivalence or grammatical quality. Repeated prefixes, incomplete definitions,
and conflicting response instructions can occur. Distinct child IDs can also
contain identical strings because there is no text-level deduplication. This
restricted operator family keeps the search easy to inspect, while
limiting conclusions about stronger LLM-generated mutations or general prompt
optimization systems.

\subsection{Evaluation scope and computational complexity}
A survivor invocation considers at most $m(m-1)$ directed candidate comparisons.
If $n$ is the typical common-sample count and $B=200$, a straightforward upper
bound for bootstrap arithmetic is $O(m^2Bn)$, in addition to maintaining
observations and sorting fronts. Cost checks, previously eliminated IDs and
minimum-sample checks reduce the comparisons that proceed to a usable decision.
The implementation is intended for small pools; no wall-clock overhead benchmark
is available, so a call-neutral gate should not be described as computationally
free.

There is also an implementation distinction between the survivor pool and the
evaluation registry. The search state accumulates IDs of generated prompts;
adding children does not remove historical IDs from that registry. The Racing
policy constructs contenders from this state, while survivor selection receives
the current retained prompts and new offspring. Thus, a removed prompt can
remain visible to subsequent evaluation allocation without being eligible to
breed or become the final winner. The present results describe that exact
implementation. Pruning evaluation eligibility to the current survivor pool
would change budget allocation and requires a new experiment, not a retrospective
reinterpretation of these logs.
